\section{Билет 2. Колебательное движение. Гармонические колебания. Горизонтальный пружинный маятник. Уравнение гармонических колебаний (с выводом).}

\begin{definition}
\textbf{Гармоническими колебаниями} называются колебания, при которых колеблющаяся величина изменяется со временем по закону синуса или косинуса. Простейшим примером системы, где возникают свободные гармонические колебания, является движение тела под действием силы упругости.
\end{definition}

\begin{theorem}[Уравнение движения горизонтального пружинного маятника]
Рассмотрим груз массы $m$, закреплённый на горизонтальной пружине жёсткостью $k$, колеблющийся по гладкой поверхности (трением пренебрегаем). Возвращающая сила описывается законом Гука:
\begin{equation}
    F_x = -kx,
\end{equation}
где знак «минус» означает, что сила всегда направлена к положению равновесия.
\begin{itemize}
    \item Если $x > 0$, то $F < 0$ и ускорение $w < 0$.
    \item Если $x < 0$, то $F > 0$ и ускорение $w > 0$.
    \item В положении равновесия $x = 0$, сила и ускорение равны нулю.
\end{itemize}

Запишем второй закон Ньютона в проекции на ось $x$:
\begin{equation}
    m w = F_x \quad \Rightarrow \quad m \frac{d^2x}{dt^2} = -kx.
\end{equation}
Разделив на $m$ и перенеся все члены в одну сторону, получаем дифференциальное уравнение:
\begin{equation}
    \frac{d^2x}{dt^2} + \frac{k}{m}x = 0.
\end{equation}
Вводя обозначение $\omega_0^2 = \frac{k}{m}$, приходим к каноническому виду уравнения гармонических колебаний:
\begin{equation}
    \ddot{x} + \omega_0^2 x = 0. \label{eq:harmonic_horiz}
\end{equation}
\end{theorem}

\begin{remark}
Величина $\omega_0 = \sqrt{k/m}$ называется \textbf{циклической (круговой) частотой} собственных колебаний. Силы вида $F_x = -kx$, имеющие неупругую природу, называют \textbf{квазиупругими}. Уравнение \eqref{eq:harmonic_horiz} универсально для всех систем с квазиупругой возвращающей силой.
\end{remark}
